\section*{Introduction}

Continuous quantum measurement has emerged as a powerful route to engineer nonclassical states and to extract metrologically useful information from noisy platforms. Rather than treating measurement backaction as a nuisance, one can tailor the measurement channel and its unravelling to steer a system toward entanglement and spin squeezing. This thesis explores that idea in the dephasing-dominated regime, focusing on \emph{entangling unravellings} that convert which-path erasure into a resource for quantum state engineering and quantum metrology.

We study a minimal, experimentally motivated setting: two qubits that radiate into modes subsequently mixed on a balanced beam splitter and measured with different continuous readouts. The only decoherence is single-qubit dephasing (local Pauli-$z$ noise). As a concrete baseline we initialize the system in the separable state \(|++\rangle\) and then generalize to coherent spin states (CSS) at varying polar angles. We compare three measurement schemes—photodetection (quantum jumps), homodyne (single-quadrature), and heterodyne (double-quadrature)—which realize distinct stochastic unravellings of the same master equation and probe symmetric/antisymmetric combinations of the local dephasing channels. In doing so, we highlight the information–backaction trade-off between homodyne and heterodyne detection and show how interferometric mixing can erase which-path information and promote conditional entanglement.

Methodologically, we formulate the stochastic master equation (SME) with an explicit split between monitored and unmonitored channels, allowing us to model finite detection efficiency as a single parameter without changing the underlying system dynamics. Analytical results are derived in the two-qubit case: we obtain closed-form jump probabilities and conditional states for the interferometric photodetection channels, and we provide compact expressions and bounds for trajectory-averaged entanglement and spin squeezing. These predictions are validated against numerical quantum-trajectory simulations implemented in \textsf{QuTiP}, ensuring internal consistency across jump and diffusive unravellings and quantifying robustness to inefficiency.

From the metrological perspective, we track standard figures of merit such as concurrence and the Kitagawa–Ueda and Wineland spin-squeezing parameters, always evaluating them in the instantaneous mean-spin frame to capture the rotation of the optimal squeezing axis induced by measurement backaction. Moving beyond the specific \(|++\rangle\) input, we map how the entanglement yield and squeezing depth depend on the CSS orientation and on the choice of readout. This reveals simple design rules for maximizing performance under dephasing: the alignment of the input state relative to the measured quadrature, the interferometric phase that selects symmetric versus antisymmetric channels, and minimum efficiency thresholds required to overcome information loss.

Finally, we outline a scalable interferometric architecture for atomic ensembles that implements a collective \(J_z\) measurement together with complementary observables accessible via alternative unravellings. We extend our analysis beyond two qubits to assess concurrence and conditional squeezing in larger systems, discussing scaling with the number of atoms and the role of measurement inefficiency in the many-body limit. The overall picture is that, even with pure dephasing as the dominant noise, appropriately engineered continuous measurements can serve as a practical and platform-agnostic resource for quantum state preparation and precision sensing.

\paragraph{Contributions.}
\begin{itemize}
	\item A unified, interferometric framework to compare photodetection, homodyne, and heterodyne \emph{entangling unravellings} under single-qubit dephasing.
	\item Closed-form analytics for two qubits (jump probabilities, conditional states) with quantitative, trajectory-averaged benchmarks for concurrence and spin squeezing.
	\item An explicit efficiency model via monitored/unmonitored channel splitting in the SME, with validated thresholds from \textsf{QuTiP} simulations.
	\item Generalization from the \(|++\rangle\) input to CSS with varying polar angle, identifying orientation-dependent optima and readout-specific trade-offs.
	\item Extension to \(N>2\) qubits and a scalable readout that realizes collective \(J_z\) together with complementary observables relevant for metrology.
\end{itemize}

\paragraph{Structure of the thesis.}
Chapter~1 reviews continuous measurement theory and stochastic unravellings. Chapter~2 introduces the two-qubit interferometric model and derives analytic results under dephasing. Chapter~3 presents numerical validations and the treatment of finite efficiency. Chapter~4 generalizes to CSS inputs and discusses metrological performance. Chapter~5 extends to multi-qubit ensembles and outlines scalable architectures and applications.
